\section{Auswertung der Usability-Tests}
\label{sec:BB-07_usability-analysis}

Dieses Kapitel dokumentiert die aggregierte Auswertung der Usability-Tests und stellt die rechnerische Herleitung der in \textbf{\hyperref[ch:05_analysis-and-results]{Kapitel~5}} berichteten Kennzahlen dar. Ziel ist die transparente Nachvollziehbarkeit der quantitativen Befunde auf Basis der Rohdaten.

% -----------------------------------------------------
\subsection{Methodisches Vorgehen}

Alle Teilnehmenden bearbeiteten fünf Aufgaben im alten sowie im neuen Wiki (n = 10, insgesamt 100 Messungen). Erfasst wurden:
\begin{itemize}
  \item Bearbeitungszeit in Sekunden,
  \item Anzahl der Fehler,
  \item Anzahl der Klicks,
  \item subjektive Schwierigkeit (Skala 1--5).
\end{itemize}

Für jede Aufgabe wurden arithmetische Mittelwerte, Mediane, Minimal- und Maximalwerte sowie Standardabweichungen berechnet. Zusätzlich wurden Gesamtfehler, durchschnittliche Fehler pro Person sowie relative Verbesserungen zwischen alter und neuer Version bestimmt. Die Auswertung ist deskriptiv angelegt und dient der strukturierten Gegenüberstellung der Nutzungseffizienz und -effektivität.

% -----------------------------------------------------
\subsection{Gesamtvergleich Alt vs. Neu}

\begin{table}[H]
\centering
\caption{Aggregierte Usability-Kennzahlen ueber alle Aufgaben}
\label{tab:BB-usability-overall}
\begin{tabular}{l c c}
\toprule
Kennzahl & Altes Wiki & Neues Wiki \\
\midrule
Mittlere Bearbeitungszeit (s) & 60,68 & 19,18 \\
Median Bearbeitungszeit (s) & 48 & 17 \\
Gesamtfehler & 100 & 15 \\
Durchschnittliche Klicks & 5,04 & 1,16 \\
\bottomrule
\end{tabular}
\end{table}

Die mittlere Bearbeitungszeit reduzierte sich um rund 68~\%, die Fehlerzahl um rund 85~\%. Auch die durchschnittliche Klickzahl verringerte sich deutlich, was auf effizientere Navigationspfade im neuen Wiki hindeutet.

% -----------------------------------------------------
\subsection{Aufgabenbezogene Zeit- und Fehlerkennzahlen}

% ---------- Aufgabe 1 ----------
\begin{table}[H]
\centering
\caption{Aufgabe 1 -- Zeit- und Fehlerkennzahlen}
\label{tab:BB-task1-stats}
\begin{tabular}{l c c}
\toprule
Kennzahl & Altes Wiki & Neues Wiki \\
\midrule
Arithmetisches Mittel (s) & 19,5 & 11,5 \\
Median (s) & 13,0 & 8,5 \\
Minimum (s) & 8 & 5 \\
Maximum (s) & 45 & 28 \\
Standardabweichung (s) & 14,28 & 7,38 \\
Gesamtfehler & 3 & 1 \\
Durchschnitt Fehler pro Person & 0,3 & 0,1 \\
Relative Zeitverbesserung & -- & -41,0\% \\
\bottomrule
\end{tabular}
\end{table}

% ---------- Aufgabe 2 ----------
\begin{table}[H]
\centering
\caption{Aufgabe 2 -- Zeit- und Fehlerkennzahlen}
\label{tab:BB-task2-stats}
\begin{tabular}{l c c}
\toprule
Kennzahl & Altes Wiki & Neues Wiki \\
\midrule
Arithmetisches Mittel (s) & 76,3 & 23,6 \\
Median (s) & 72,5 & 22,0 \\
Minimum (s) & 46 & 13 \\
Maximum (s) & 138 & 34 \\
Standardabweichung (s) & 28,99 & 7,86 \\
Gesamtfehler & 20 & 3 \\
Durchschnitt Fehler pro Person & 2,0 & 0,3 \\
Relative Zeitverbesserung & -- & -69,1\% \\
\bottomrule
\end{tabular}
\end{table}

% ---------- Aufgabe 3 ----------
\begin{table}[H]
\centering
\caption{Aufgabe 3 -- Zeit- und Fehlerkennzahlen}
\label{tab:BB-task3-stats}
\begin{tabular}{l c c}
\toprule
Kennzahl & Altes Wiki & Neues Wiki \\
\midrule
Arithmetisches Mittel (s) & 132,4 & 13,8 \\
Median (s) & 132,0 & 12,5 \\
Minimum (s) & 48 & 6 \\
Maximum (s) & 180 & 28 \\
Standardabweichung (s) & 51,07 & 7,10 \\
Gesamtfehler & 51 & 0 \\
Durchschnitt Fehler pro Person & 5,1 & 0,0 \\
Relative Zeitverbesserung & -- & -89,6\% \\
\bottomrule
\end{tabular}
\end{table}

% ---------- Aufgabe 4 ----------
\begin{table}[H]
\centering
\caption{Aufgabe 4 -- Zeit- und Fehlerkennzahlen}
\label{tab:BB-task4-stats}
\begin{tabular}{l c c}
\toprule
Kennzahl & Altes Wiki & Neues Wiki \\
\midrule
Arithmetisches Mittel (s) & 41,9 & 16,4 \\
Median (s) & 43,0 & 16,5 \\
Minimum (s) & 16 & 5 \\
Maximum (s) & 78 & 35 \\
Standardabweichung (s) & 20,84 & 8,67 \\
Gesamtfehler & 16 & 2 \\
Durchschnitt Fehler pro Person & 1,6 & 0,2 \\
Relative Zeitverbesserung & -- & -60,9\% \\
\bottomrule
\end{tabular}
\end{table}

% ---------- Aufgabe 5 ----------
\begin{table}[H]
\centering
\caption{Aufgabe 5 -- Zeit- und Fehlerkennzahlen}
\label{tab:BB-task5-stats}
\begin{tabular}{l c c}
\toprule
Kennzahl & Altes Wiki & Neues Wiki \\
\midrule
Arithmetisches Mittel (s) & 33,3 & 30,6 \\
Median (s) & 29,0 & 26,0 \\
Minimum (s) & 18 & 13 \\
Maximum (s) & 58 & 58 \\
Standardabweichung (s) & 12,91 & 14,51 \\
Gesamtfehler & 10 & 9 \\
Durchschnitt Fehler pro Person & 1,0 & 0,9 \\
Relative Zeitverbesserung & -- & -8,1\% \\
\bottomrule
\end{tabular}
\end{table}

% -----------------------------------------------------
\newpage
\subsection{Klickkennzahlen je Aufgabe}

% ---------- Klicks Aufgabe 1 ----------
\begin{table}[H]
\centering
\caption{Aufgabe 1 -- Klickkennzahlen}
\label{tab:BB-task1-clicks}
\begin{tabular}{l c c}
\toprule
Kennzahl & Altes Wiki & Neues Wiki \\
\midrule
Arithmetisches Mittel & 1,9 & 0,5 \\
Median & 1,5 & 0,0 \\
Minimum & 0 & 0 \\
Maximum & 8 & 3 \\
Standardabweichung & 2,28 & 0,97 \\
Relative Klickreduktion & -- & -73,7\% \\
\bottomrule
\end{tabular}
\end{table}

% ---------- Klicks Aufgabe 2 ----------
\begin{table}[H]
\centering
\caption{Aufgabe 2 -- Klickkennzahlen}
\label{tab:BB-task2-clicks}
\begin{tabular}{l c c}
\toprule
Kennzahl & Altes Wiki & Neues Wiki \\
\midrule
Arithmetisches Mittel & 5,2 & 1,6 \\
Median & 5,0 & 1,0 \\
Minimum & 2 & 0 \\
Maximum & 10 & 5 \\
Standardabweichung & 2,35 & 1,35 \\
Relative Klickreduktion & -- & -69,2\% \\
\bottomrule
\end{tabular}
\end{table}

% ---------- Klicks Aufgabe 3 ----------
\begin{table}[H]
\centering
\caption{Aufgabe 3 -- Klickkennzahlen}
\label{tab:BB-task3-clicks}
\begin{tabular}{l c c}
\toprule
Kennzahl & Altes Wiki & Neues Wiki \\
\midrule
Arithmetisches Mittel & 9,0 & 0,3 \\
Median & 11,0 & 0,0 \\
Minimum & 2 & 0 \\
Maximum & 18 & 1 \\
Standardabweichung & 5,25 & 0,48 \\
Relative Klickreduktion & -- & -96,7\% \\
\bottomrule
\end{tabular}
\end{table}

% ---------- Klicks Aufgabe 4 ----------
\begin{table}[H]
\centering
\caption{Aufgabe 4 -- Klickkennzahlen}
\label{tab:BB-task4-clicks}
\begin{tabular}{l c c}
\toprule
Kennzahl & Altes Wiki & Neues Wiki \\
\midrule
Arithmetisches Mittel & 4,9 & 1,3 \\
Median & 3,5 & 1,0 \\
Minimum & 1 & 1 \\
Maximum & 10 & 3 \\
Standardabweichung & 3,35 & 0,67 \\
Relative Klickreduktion & -- & -73,5\% \\
\bottomrule
\end{tabular}
\end{table}

% ---------- Klicks Aufgabe 5 ----------
\begin{table}[H]
\centering
\caption{Aufgabe 5 -- Klickkennzahlen}
\label{tab:BB-task5-clicks}
\begin{tabular}{l c c}
\toprule
Kennzahl & Altes Wiki & Neues Wiki \\
\midrule
Arithmetisches Mittel & 4,2 & 2,1 \\
Median & 3,0 & 2,0 \\
Minimum & 2 & 1 \\
Maximum & 9 & 4 \\
Standardabweichung & 2,66 & 0,88 \\
Relative Klickreduktion & -- & -50,0\% \\
\bottomrule
\end{tabular}
\end{table}

% -----------------------------------------------------
\subsection{Zusammenfassende Einordnung}

Die Auswertung zeigt über alle Aufgaben hinweg deutliche Verbesserungen der Nutzungseffizienz und -effektivität im neuen Wiki. Besonders ausgeprägt sind die Effekte bei Aufgaben mit hohem Such- und Orientierungsaufwand (Aufgaben~2 und~3). Aufgabe~5 weist dagegen nur geringe Verbesserungen auf und stellt einen stabilen Ausreißer dar. Dieses Muster korrespondiert mit den qualitativen Aussagen zur eingeschränkten Verständlichkeit der Personen- und Ansprechpartnerlogik.

\newpage
